\section{Round-sixteen positive closure: pre-contact transversality and positivity-preserving recovery}
\label{sec:r16-b2}

The first-surplus geometry is formulated before the surplus collision occurs.
On an open parameter domain the two candidate particles follow genuine
hard-sphere trajectories only through the prescribed earlier contacts.  The
surplus contact is imposed as a zero set; reflection and all later dynamics are
introduced only on that set.  Lower recovery is an entropy projection onto the
positive balanced cone, never an unconstrained signed correction.

\subsection{Compatible Green traces}

For each particle number let \(\Omega_{\varepsilon,N}\) be the hard-core phase
domain.  On a binary boundary face \(\Sigma_{ij}\), write
\(n_{ij}=(x_i-x_j)/\varepsilon\) and
\(g_{ij}=(v_i-v_j)\cdot n_{ij}\).  The incoming trace measure is the normal
flux
\[
 d\gamma^-_{ij}=(-g_{ij})_+ f_N\,d\sigma_{ij}
\]
and is not multiplied by a second flux factor.  Incoming and outgoing traces
are related by the specular involution \(R_{ij}\).

\begin{theorem}[Closed hard-sphere Green graph]
\label{thm:r16-b2-trace}
On the weighted \(L^1\) Liouville space, the free transport operator with
specular boundary condition has a closed graph domain whose normal traces are
finite flux measures and satisfy
\[
 \int_{\Omega_N} (V\cdot\nabla_X\phi)f_N
 +\int_{\Omega_N}\phi(V\cdot\nabla_X f_N)
 =\sum_{i<j}\int_{\Sigma_{ij}}(\phi^+-\phi^-)\,d\gamma^-_{ij}.
\]
The graph is stable under the bounded path/contact tilts used below.  Multiple
and grazing faces are retained as graph strata; finite entropy contact flows
are absolutely continuous with respect to the simple incoming flux and hence
charge no higher-codimension stratum.
\end{theorem}

\begin{proof}
Apply the divergence-measure Green formula on smooth truncations which exclude
an \(\eta\)-neighborhood of corners and grazing.  The boundary term is the
normal trace of \(Vf_N\), hence already includes \(|g_{ij}|\).  Specular
reflection preserves surface measure, kinetic energy, and the absolute flux.
Uniform weighted moment bounds let \(\eta\downarrow0\); codimension-two strata
have zero reference flux.  Closure follows from convergence of the interior
divergence and trace measures in the Green identity.  A bounded exponential
source changes the graph norm by an equivalent weight.  If a limiting contact
measure had mass on a stratum where the reference simple flux vanishes, its
relative entropy would be infinite by singular-part duality.
\end{proof}

\subsection{Open pre-contact parameter spaces}

Fix a labelled creation genealogy \(\mathcal G\) and a chronological first
surplus pair \((a,b)\) at candidate time \(t_*\).  Let \(\Theta_{\mathcal G}\)
consist of root positions and velocities, creation times, incoming normals, and
impact parameters for all contacts strictly before \(t_*\), with strict time
ordering, non-grazing inequalities, and clearance from every unintended
contact.  It is open.  Propagate the true hard-sphere flow only through those
earlier contacts.  Define the pre-contact map
\[
 \mathfrak F_{\mathcal G}(\theta,t)
 =x_a^-(t;\theta)-x_b^-(t;\theta),
 \qquad (\theta,t)\in\Theta_{\mathcal G}\times I_*.
\]
It is smooth on the open domain.  The surplus contact manifold is
\[
 \mathcal C_{\mathcal G}
 =\{(\theta,t):|\mathfrak F_{\mathcal G}(\theta,t)|^2
   =\varepsilon^2,
   \mathfrak F_{\mathcal G}(\theta,t)
    \cdot(v_a^--v_b^-)<0\}.
\]

\begin{lemma}[Physical pre-contact submersion]
\label{lem:r16-b2-submersion}
For every finite regular genealogy there is a finite cover of
\(\Theta_{\mathcal G}\) by charts and, on each chart, a pair of genuine
pre-contact variables \(\zeta\in\mathbb R^2\) such that on
\(\mathcal C_{\mathcal G}\)
\[
 \left|\det D_\zeta
 \Pi_{(v_a^--v_b^-)^\perp}
 \mathfrak F_{\mathcal G}(\theta,t)\right|
 \ge c_{\mathcal G}>0
\]
outside a set whose incoming-flux integral is at most
\(C_{\mathcal G}\delta^{\alpha_{\mathcal G}}\).  Reflection and the future
trajectory are not used in defining this derivative.
\end{lemma}

\begin{proof}
Differentiate the true earlier collision chronology by the hard-sphere Jacobi
maps.  Root translations, root velocities, and the two tangential coordinates
at the youngest separating creation fork form the available variables.  If
all transverse two-minors vanished identically, the exterior square of every
available Jacobi column would vanish.  Propagating the annihilating covector
backward through invertible free-flight and non-grazing reflection maps would
make it annihilate two independent root translation/velocity directions, a
contradiction.  Hence one analytic minor is not identically zero.

For a concrete witness, choose root particles in disjoint narrow tubes, assign
strictly ordered rational flight times, and solve each creation collision
successively by the implicit reflection equation.  At every step choose the
incoming normal from an open cone and shorten the next tube so all unlisted
pairs remain separated by \(2\varepsilon\).  The final two free tangential
variables move particles \(a,b\) independently in their relative transverse
plane, giving determinant one at the zero-interaction center.  Openness
provides a genuine regular hard-sphere chronology.  Analytic sublevel
estimates on the compact chart then give the stated exceptional flux bound.
Only after restricting to \(\mathcal C_{\mathcal G}\) do we set
\(n=\mathfrak F/\varepsilon\), apply the specular reflection, and use the
implicit-function theorem to continue every later prescribed contact.
\end{proof}

\begin{lemma}[First-surplus coarea gain]
\label{lem:r16-b2-gain}
For a fixed genealogy, energy cutoff \(\mathbb V\), and bounded contact
source, integration over the first surplus contact gives
\[
 C_{\mathcal G}(1+\mathbb V)^{C_{\mathcal G}}
 (1+T)^{C_{\mathcal G}}\varepsilon^{\alpha_{\mathcal G}}
\]
relative to the creation-tree measure.  The estimate retains the true
post-contact reflection and every later actual contact on the contact zero set.
\end{lemma}

\begin{proof}
On a chart from Lemma~\ref{lem:r16-b2-submersion}, apply coarea to the two
transverse components of the pre-contact map and the scalar contact time.  The
incoming normal flux cancels the time-Jacobian degeneracy.  The regular region
gives the tubular \(\varepsilon\)-gain; the analytic-minor sublevel set gives
\(\varepsilon^{\alpha_{\mathcal G}}\) after optimizing its cutoff.  On the
zero set the reflection map is smooth and the later chronology is the
implicit-function continuation constructed above, so the integrand is the
true reflected future.  No off-contact reflection is evaluated.
\end{proof}

\subsection{Fixed-horizon dominated genealogy summation}

Let \(w(\mathcal G)\) be the absolute creation-tree activity including all
ordered time simplices, Gaussian velocity weights, homogeneity cutoffs, and
source marks.

\begin{theorem}[All-contact fixed-horizon majorant]
\label{thm:r16-b2-majorant}
For \(T\le T_0\),
\[
 \sum_{\mathcal G}w(\mathcal G)<\infty,
 \qquad
 \sum_{\mathcal G}w(\mathcal G)
 C_{\mathcal G}(1+\mathbb V)^{C_{\mathcal G}}
 \varepsilon^{\alpha_{\mathcal G}}\longrightarrow0.
\]
The limits remain valid after removing velocity and grazing cutoffs and after
any fixed number of path/contact source derivatives.  Consequently the sum of
all cyclic genealogies is \(o(1)\) on the full fixed horizon.
\end{theorem}

\begin{proof}
Each creation contact contributes a collision coarea operator bounded by
\(CT\) in the Gaussian velocity norm; the ordered time simplex contributes
\(1/m!\).  Tree enumeration gives
\(w(\mathcal G)\le k!C^kT^{k-1}/m!\).  Incoming flux controls grazing because
the singular inverse collision derivative is paired with its normal velocity
factor.  Gaussian moments control every polynomial velocity loss.  Thus one
summable sequence dominates all cutoff integrands and all fixed source
derivatives.  For each fixed genealogy, Lemma~\ref{lem:r16-b2-gain} tends to
zero.  Dominated convergence over the countable genealogy family proves the
second assertion, then monotone exhaustion of positive absolute majorants
removes the cutoffs.  Complex activities are handled only after taking this
absolute majorant, not by a signed monotone-convergence argument.
\end{proof}

\begin{theorem}[Grand-canonical actual-contact pressure]
\label{thm:r16-b2-pressure}
The grand-canonical joint density/actual-contact log-Laplace functional has a
normally convergent fixed-horizon expansion on a common complex source ball.
It converges with all fixed source derivatives to the unique analytic solution
of the Hamilton--Jacobi equation with Hamiltonian
\[
 \mathbb H(f,p,\psi)=
 \langle f,v\cdot\nabla_xp\rangle
 +{1\over2}\int ff_*B
 \bigl(e^{\Delta p+\psi}-1\bigr).
\]
Every \(\psi\)-derivative selects an actual incoming contact.
\end{theorem}

\begin{proof}
Connected trajectory clusters are organized by creation genealogies and
chronological surplus edges.  Creation trees converge by collision coarea.
The first surplus edge of every cyclic cluster is controlled by
Theorem~\ref{thm:r16-b2-majorant}, so all cyclic contributions vanish after
any fixed source differentiation.  Removing the last collision from a
limiting tree leaves two independent rooted subtrees and gives the displayed
Hamiltonian.  Normal convergence permits differentiation and Picard uniqueness
on the analytic source scale.
\end{proof}

\subsection{Positive balanced recovery by entropy projection}

For a positive density path \(f\), let
\[
 dA_f={1\over2}ff_*B\,dt\,dx\,dv\,dv_*\,d\omega.
\]
Let \(\mathcal B(f,\Gamma)=0\) denote the complete weak balance relation,
including endpoints.  Starting from a smooth strictly positive approximate
intensity \(q_0A_f\), minimize
\[
 \mathcal J(q)=\int\ell(q)\,dA_f,
 \qquad \ell(q)=q\log q-q+1,
\]
over \(q\ge0\) subject to \(\mathcal B(f,qA_f)=0\).

\begin{theorem}[Positive entropic balance projection]
\label{thm:r16-b2-projection}
On the regular short-time phase chart, the feasible positive balanced set is
nonempty and the entropy projection has a unique minimizer
\[
 q^*=q_0\exp\{\Delta r^*\}.
\]
The multiplier \(r^*\) lies in the closed balance quotient and satisfies a
nonlinear coercive Euler equation.  If the balance defect of \(q_0A_f\) tends
to zero, then \(q^*A_f-q_0A_f\to0\) in weighted variation and
\(\mathcal J(q^*)-\mathcal J(q_0)\to0\).  Positivity is automatic and no
pointwise bound on an additive correction is required.
\end{theorem}

\begin{proof}
The functional is strictly convex, lower semicontinuous, and superlinear on
\(L^1(A_f)\).  The balance constraint is closed and affine.  A positive
reference feasible intensity is obtained by a small Maxwellian collision bath
plus the corresponding transport solution, so the direct method gives a
unique minimizer.  Fenchel duality yields
\(\log(q^*/q_0)=\Delta r^*\).  The finite-time collision/transport
observability estimate on the regular phase chart makes the dual functional
strictly coercive modulo collision invariants, giving existence and uniqueness
of the multiplier class.  Relative-entropy Pinsker control and the dual
Euler equation bound the distance to \(q_0\) by the balance defect.  Because
the correction is multiplicative, \(q^*>0\) regardless of the size of
\(\Delta r^*\).
\end{proof}

\begin{theorem}[Source exposure and entropy exhaustion]
\label{thm:r16-b2-exposure}
Every smooth balanced pair with
\(0<c\le q=d\Gamma/dA_f\le C\) and regular coefficients is generated by a
bounded contact source
\[
 \psi=\log q-\Delta p
\]
for a solution of the backward adjoint equation.  Every finite-action balanced
pair is approximated by such exposed pairs with convergence of state and
action.  The approximations use velocity truncation, positive mollification,
and Theorem~\ref{thm:r16-b2-projection}.  Hence increasing bounded source
balls exhaust the full entropy rate.
\end{theorem}

\begin{proof}
For a regular positive pair, solve the nonautonomous backward adjoint on the
weighted graph domain; short-time coercivity gives bounded \(p\) and
\(\Delta p\), so \(\psi\) is bounded.  The first derivative of
Theorem~\ref{thm:r16-b2-pressure} is then the prescribed controlled Boltzmann
pair by uniqueness of the biased forward equation.

For a finite-action pair, truncate velocities and \(q\), convolve away from
endpoints, and add a vanishing positive Maxwellian background.  This creates a
small balance defect while preserving entropy convergence.  Apply the positive
entropy projection to restore balance exactly.  Pinsker control makes the
projection error vanish, and lower semicontinuity plus the recovery upper
bound gives convergence of the action.  The required source radii increase
with the truncation, yielding monotone dual exhaustion.
\end{proof}

\subsection{The joint LDP}

\begin{theorem}[Grand-canonical density--contact LDP]
\label{thm:r16-b2-gc-ldp}
The pair \((\pi^\varepsilon,\Gamma^\varepsilon)\) satisfies a full good LDP on
the weighted graph state space with speed \(\mu_\varepsilon\) and rate
\[
 I_T^{\rm gc}(f,\Gamma)
 =I_0^{\rm gc}(f_0)+
 \begin{cases}
 \displaystyle\int\ell(d\Gamma/dA_f)dA_f,
 &\mathcal B(f,\Gamma)=0,
 \ \Gamma\ll A_f,\\
 +\infty,&\text{otherwise}.
 \end{cases}
\]
Finite action excludes grazing, corner, and multiple-contact singular mass.
\end{theorem}

\begin{proof}
The analytic pressure and exponential moment sources give exponential
compactness and the upper bound.  Integral Fenchel duality imposes balance and
produces the entropy integrand.  The exposed-pair change of measure gives the
lower bound on the regular class; under the tilted law, normal convergence of
the source expansion gives concentration.  Theorem~\ref{thm:r16-b2-exposure}
extends the lower bound to every finite-action pair.  Singular contact strata
are outside the support of \(A_f\), so any mass there has infinite singular
entropy.
\end{proof}

\begin{theorem}[Source-conditioned microcanonical LDP]
\label{thm:r16-b2-mc-ldp}
After the B1 exact-number source-dependent conditioning theorem, the same
joint LDP holds for the primitive microcanonical shell, with the initial term
replaced by the constrained initial rate.  The proof order is
\[
 \text{B2-GC}\longrightarrow\text{B1}\longrightarrow\text{B2-MC}.
\]
\end{theorem}

\begin{proof}
B1 supplies local-uniform conditioned log-Laplace convergence for every
particle/contact source and a subexponential centered shell coefficient under
the reoptimized tilted law.  Apply it to the upper-bound sources and to each
regular lower-bound exposing source.  The dynamic entropy action is unchanged;
only the initial rate is contracted by the shell constraints.  Source
exhaustion and positive recovery pass through the same conditioning.
\end{proof}

\begin{theorem}[Round-sixteen B2 closure]
\label{thm:r16-b2-main}
Actual contacts are controlled through a genuine open pre-contact evaluation
map, reflection is introduced only on its contact zero set, fixed genealogies
are summed by a cutoff-uniform factorial majorant, and all finite-action lower
recoveries are positive entropy projections.  This closes the grand-canonical
root theorem and its noncircular microcanonical transfer.
\end{theorem}

\begin{proof}
Combine Theorems~\ref{thm:r16-b2-trace},
\ref{thm:r16-b2-majorant}, \ref{thm:r16-b2-pressure},
\ref{thm:r16-b2-projection}, \ref{thm:r16-b2-exposure},
\ref{thm:r16-b2-gc-ldp}, and \ref{thm:r16-b2-mc-ldp}, and
Lemmas~\ref{lem:r16-b2-submersion} and \ref{lem:r16-b2-gain}.
\end{proof}
